\section{Decyzje architektoniczne}

\subsection{Podział na warstwy}
\begin{figure}[H]
\centering
\includegraphics[width=0.8\textwidth]{wybor_architektury.png}
\caption{Diagram podziału na warstwy}
\label{fig:layers}
\end{figure}

\begin{itemize}
\item Warstwa prezentacji: trzy różne aplikacje klienckie — aplikacja webowa, mobilna i aplikacja na urządzenia embedded.
\item Warstwa logiki biznesowej: centralna aplikacja serwerowa (serverApp), w której znajdują się komponenty takie jak UserService, PaymentService, PositionService itd.
\item Warstwa dostępu do danych: komponenty UserRepository, PaymentRepository.
\item Warstwa danych: baza danych Oracle z mechanizmem failover.
\end{itemize}

\textbf{Decyzja:} Rozdzielenie odpowiedzialności na warstwy zwiększa modularność i ułatwia zarządzanie kodem oraz jego testowanie.

\textbf{Alternatywa:} Architektura mikroserwisowa

\textbf{Zalety alternatywy:}
\begin{itemize}
\item Lepsza skalowalność poszczególnych komponentów
\item Możliwość niezależnego wdrażania i rozwijania usług
\item Lepsza odporność na awarie (awaria jednego serwisu $\neq$ awaria całego systemu)
\end{itemize}

\textbf{Wady alternatywy:}
\begin{itemize}
\item Większa złożoność wdrożeniowa (DevOps, CI/CD, monitoring)
\item Konieczność rozwiązywania problemów związanych z komunikacją między serwisami
\item Trudniejsze debugowanie i testowanie end-to-end
\end{itemize}

\subsection{Modularność i komponenty (Component-based Design)}
\begin{figure}[H]
\centering
\includegraphics[width=0.8\textwidth]{organizacja_aplikacji_serwerowej.png}
\caption{Diagram komponentów}
\label{fig:components}
\end{figure}

\begin{itemize}
\item Serwerowa aplikacja została podzielona na komponenty pełniące konkretne role (SigninController, TollController, MainComponent, itd.).
\item Każdy komponent ma jasno zdefiniowaną odpowiedzialność, zgodnie z zasadą Single Responsibility Principle.
\end{itemize}

\textbf{Decyzja:} Wprowadzenie komponentów umożliwia łatwe rozszerzanie i testowanie poszczególnych fragmentów systemu.

\textbf{Alternatywa:} Monolityczna aplikacja serwerowa

\textbf{Zalety alternatywy:}
\begin{itemize}
\item Prostsza implementacja i wdrożenie
\item Mniejsza liczba zależności i konfiguracji
\item Mniej złożone środowisko developerskie
\end{itemize}

\textbf{Wady alternatywy:}
\begin{itemize}
\item Trudniejsza skalowalność i refaktoryzacja
\item Każda zmiana wymaga redeploy całej aplikacji
\item Trudniejsze testowanie izolowanych funkcji
\end{itemize}

\subsection{Wielokanałowy dostęp (Multi-Platform Clients)}
\begin{figure}[H]
\centering
\includegraphics[width=0.8\textwidth]{dostep_do_systemu.png}
\caption{Diagram wielokanałowego dostępu}
\label{fig:multichannel}
\end{figure}

\begin{itemize}
\item Użytkownicy mogą korzystać z systemu za pomocą aplikacji mobilnej, webowej lub urządzeń embedded.
\end{itemize}

\textbf{Decyzja:} Umożliwienie różnym grupom użytkowników (np. administratorzy vs kierowcy) dostępu do funkcji systemu w najbardziej dogodny sposób.

\textbf{Alternatywa:} Tylko aplikacja mobilna (np. PWA lub natywna)

\textbf{Zalety alternatywy:}
\begin{itemize}
\item Uproszczony interfejs użytkownika
\item Jedna platforma do utrzymania
\item Lepsze dopasowanie do kontekstu użytkownika (kierowcy)
\end{itemize}

\textbf{Wady alternatywy:}
\begin{itemize}
\item Brak wygodnego interfejsu dla administratorów lub analityków
\item Mniejsza elastyczność użytkowania
\item Trudności z dostępem do systemu z urządzeń stacjonarnych
\end{itemize}

\subsection{Rozdzielenie ról i uprawnień}
\begin{figure}[H]
\centering
\includegraphics[width=0.8\textwidth]{uprawnienia_dla_uzytkownikow.png}
\caption{Diagram ról i uprawnień}
\label{fig:roles}
\end{figure}

\begin{itemize}
\item Administrator ma dostęp tylko przez aplikację webową.
\item Kierowca może korzystać z trzech różnych interfejsów — w zależności od potrzeb.
\end{itemize}

\textbf{Decyzja:} Jasny podział ról zwiększa bezpieczeństwo i ergonomię systemu.

\textbf{Alternatywa:} Jeden zunifikowany interfejs z uprawnieniami na poziomie konta

\textbf{Zalety alternatywy:}
\begin{itemize}
\item Mniejsze zróżnicowanie UI
\item Mniej kodu i testów związanych z różnymi platformami
\end{itemize}

\textbf{Wady alternatywy:}
\begin{itemize}
\item Mniejsza przejrzystość
\item Możliwość przypadkowego ujawnienia funkcji nieprzeznaczonych dla danego typu użytkownika
\end{itemize}

\subsection{Integracja z zewnętrznymi systemami}
\begin{figure}[H]
\centering
\includegraphics[width=0.8\textwidth]{realizacja_obslugi_platnosci_i_archiwizacji_danych.png}
\caption{Diagram integracji z systemami zewnętrznymi}
\label{fig:integration}
\end{figure}

\begin{itemize}
\item System płatności (systemPlatnosci) oraz system archiwizacji danych (archiwum) są zewnętrznymi systemami zintegrowanymi z systemem e-Toll.
\end{itemize}

\textbf{Decyzja:} Wydzielenie tych odpowiedzialności do zewnętrznych systemów pozwala na lepsze skalowanie oraz wykorzystanie istniejących rozwiązań.

\textbf{Alternatywa:} Wszystko w ramach jednej aplikacji (np. własny moduł płatności i archiwizacji)

\textbf{Zalety alternatywy:}
\begin{itemize}
\item Większa kontrola nad logiką
\item Mniejsze zależności zewnętrzne
\end{itemize}

\textbf{Wady alternatywy:}
\begin{itemize}
\item Większe ryzyko błędów w obszarach regulowanych prawnie (np. płatności)
\item Większe koszty utrzymania i certyfikacji
\item Brak skalowalności i elastyczności
\end{itemize}

\subsection{Wysoka dostępność i odporność na awarie}
\begin{figure}[H]
\centering
\includegraphics[width=0.8\textwidth]{zapewnienie_dostepnosci_i_odpornosci_na_awarie.png}
\caption{Diagram wysokiej dostępności}
\label{fig:high-availability}
\end{figure}

\begin{itemize}
\item Baza danych jest replikowana (primary–secondary).
\item Oddzielne deployment node'y dla różnych typów urządzeń oraz dla aplikacji serwerowej i bazy danych.
\item Failover serwera bazy danych — Oracle - Secondary.
\end{itemize}

\textbf{Decyzja:} Architektura uwzględnia mechanizmy zapewniające ciągłość działania systemu nawet w przypadku awarii.

\textbf{Alternatywa:} Deployment w chmurze

\textbf{Zalety alternatywy:}
\begin{itemize}
\item Duża skalowalność
\item Bezpieczeństwo i redundancja
\end{itemize}

\textbf{Wady alternatywy:}
\begin{itemize}
\item Wyższe koszty
\item Przetwarzanie danych przez inne podmioty
\end{itemize}

\subsection{Wybór technologii wdrożeniowych}
\begin{figure}[H]
\centering
\includegraphics[width=0.8\textwidth]{wybor_technologii_wdrozeniowej.png}
\caption{Diagram technologii wdrożeniowych}
\label{fig:deployment-tech}
\end{figure}

\begin{itemize}
\item System hostowany na serwerach z Ubuntu 24.04 LTS.
\item Serwer aplikacji oparty o Apache Tomcat 8.x.
\item Baza danych Oracle 12c.
\end{itemize}

\textbf{Decyzja:} Użycie sprawdzonych technologii o długoterminowym wsparciu i wysokiej wydajności.

\textbf{Alternatywa:} Cloud-native stack – np. AWS/GCP

\textbf{Zalety alternatywy:}
\begin{itemize}
\item Automatyczne skalowanie, monitoring, CI/CD
\item Niższe koszty utrzymania fizycznej infrastruktury
\item Łatwiejsze zarządzanie kontenerami i usługami
\end{itemize}

\textbf{Wady alternatywy:}
\begin{itemize}
\item Uzależnienie od chmury (vendor lock-in)
\item Potrzeba opanowania nowych technologii
\item Potencjalne wyższe koszty początkowe
\end{itemize}

